\FigureLayoutDeclare{img/2008/shandong_science/q20_fig1.png}{8.0cm}{1bp 7bp 11bp 11bp}

\chapter{2008年山东卷（理）}

\section{单选题}

\begin{problem}
满足 \(M \subseteq \{a_1, a_2, a_3, a_4\}\)，且 \(M \cap \{a_1, a_2, a_3\} = \{a_1, a_2\}\) 的集合 \(M\) 的个数是（\quad）

\choices
    {\(1\)}
    {\(2\)}
    {\(3\)}
    {\(4\)}
\end{problem}

\begin{answer}
B
\end{answer}

\begin{solution}
由交集条件可知，\(a_1\)，\(a_2\) 必须属于 \(M\)，而 \(a_3\) 不属于 \(M\). 元素 \(a_4\) 可以属于 \(M\)，也可以不属于 \(M\)，所以 \(M\) 有 \(\{a_1,a_2\}\) 和 \(\{a_1,a_2,a_4\}\) 两种，故选 B.
\end{solution}

\begin{problem}
设复数 \(z\) 满足 \(z + \overline{z} = 4\)，\(z \cdot \overline{z} = 8\)，则 \(\frac{\overline{z}}{z}\) 等于（\quad）

\choices
    {\(\i\)}
    {\(-\i\)}
    {\(\pm 1\)}
    {\(\pm \i\)}
\end{problem}

\begin{answer}
D
\end{answer}

\begin{solution}
设 \(z=x+y\i\)（\(x,y\in\R\)），则 \(z+\overline z=2x=4\)，从而 \(x=2\). 又 \(z\overline z=x^2+y^2=8\)，所以 \(y^2=4\)，即 \(y=2\) 或 \(y=-2\). 当 \(z=2+2\i\) 时，\(\frac{\overline z}{z}=-\i\)；当 \(z=2-2\i\) 时，\(\frac{\overline z}{z}=\i\). 因此 \(\frac{\overline z}{z}=\pm\i\)，故选 D.
\end{solution}

\begin{problem}
函数 \(y = \ln \cos x\)（\(-\frac{\pi}{2} < x < \frac{\pi}{2}\)）的图象是（\quad）

\choices
    {\choicebitmap[width=0.15\paperwidth]{img/2008/shandong_science/q3_choiceA.png}}
    {\choicebitmap[width=0.15\paperwidth]{img/2008/shandong_science/q3_choiceB.png}}
    {\choicebitmap[width=0.15\paperwidth]{img/2008/shandong_science/q3_choiceC.png}}
    {\choicebitmap[width=0.15\paperwidth]{img/2008/shandong_science/q3_choiceD.png}}
\end{problem}

\begin{answer}
A
\end{answer}

\begin{solution}
在定义域 \(\left(-\frac\pi2,\frac\pi2\right)\) 上，\(\cos x>0\)，且 \(\ln\cos x\) 是偶函数. 其导数为 \(y'=-\tan x\)，所以函数在 \(\left(-\frac\pi2,0\right)\) 上单调递增，在 \(\left(0,\frac\pi2\right)\) 上单调递减，并且 \(f(0)=0\)，当 \(x\) 趋近定义域端点时函数值趋于 \(-\infty\). 因而图象关于 \(y\) 轴对称，在原点处取得最大值，且两端向下无界，故选 A.
\end{solution}

\begin{problem}
设函数 \(f(x) = |x + 1| + |x - a|\) 的图象关于直线 \(x = 1\) 对称，则 \(a\) 的值为（\quad）

\choices
    {\(3\)}
    {\(2\)}
    {\(1\)}
    {\(-1\)}
\end{problem}

\begin{answer}
A
\end{answer}

\begin{solution}
函数 \(f(x)=|x+1|+|x-a|\) 的图象由两个折点 \(x=-1\) 和 \(x=a\) 确定，且这两个折点关于图象的对称轴对称. 因此对称轴为 \(x=\frac{-1+a}{2}\). 由题设对称轴为 \(x=1\)，得 \(\frac{a-1}{2}=1\)，所以 \(a=3\)，故选 A.
\end{solution}

\begin{problem}
已知 \(\cos \left(\alpha - \frac{\pi}{6}\right) + \sin \alpha = \frac{4}{5}\sqrt{3}\)，则 \(\sin \left(\alpha + \frac{7\pi}{6}\right)\) 的值是（\quad）

\choices
    {\(-\frac{2\sqrt{3}}{5}\)}
    {\(\frac{2\sqrt{3}}{5}\)}
    {\(-\frac{4}{5}\)}
    {\(\frac{4}{5}\)}
\end{problem}

\begin{answer}
C
\end{answer}

\begin{solution}
由和角公式，题设左端为 \(\cos\left(\alpha-\frac\pi6\right)+\sin\alpha=\frac{\sqrt3}{2}\cos\alpha+\frac32\sin\alpha=\frac{\sqrt3}{2}\left(\cos\alpha+\sqrt3\sin\alpha\right)\). 因而 \(\cos\alpha+\sqrt3\sin\alpha=\frac85\). 又 \(\sin\left(\alpha+\frac{7\pi}{6}\right)=-\frac12\cos\alpha-\frac{\sqrt3}{2}\sin\alpha=-\frac12\left(\cos\alpha+\sqrt3\sin\alpha\right)=-\frac45\)，故选 C.
\end{solution}

\begin{problem}
如图是一个几何体的三视图，根据图中数据，可得该几何体的表面积是（\quad）

\bitmapfigure[max width=0.5\linewidth]{img/2008/shandong_science/q6_fig1.png}

\choices
    {\(9\pi\)}
    {\(10\pi\)}
    {\(11\pi\)}
    {\(12\pi\)}
\end{problem}

\begin{answer}
D
\end{answer}

\begin{solution}
由俯视图为圆，且正视图和侧视图均显示上、下两部分，可知该几何体由半径为 \(1\) 的球和高为 \(3\) 的圆柱组成，球与圆柱相切. 因而球的表面积为 \(4\pi\)，圆柱的两个底面面积为 \(2\pi\)，侧面积为 \(2\pi\cdot1\cdot3=6\pi\). 所以该几何体的表面积为 \(4\pi+2\pi+6\pi=12\pi\)，故选 D.
\end{solution}

\begin{problem}
在某地的奥运火炬传递活动中，有编号为 \(1\)，\(2\)，\(3\)，…，\(18\) 的 \(18\) 名火炬手. 若从中任选 \(3\) 人，则选出的火炬手的编号能组成公差为 \(3\) 的等差数列的概率为（\quad）

\choices
    {\(\frac{1}{81}\)}
    {\(\frac{1}{68}\)}
    {\(\frac{1}{306}\)}
    {\(\frac{1}{408}\)}
\end{problem}

\begin{answer}
B
\end{answer}

\begin{solution}
从 \(1\) 至 \(18\) 中任选 \(3\) 人的基本事件总数为 \(\mathrm C_{18}^{3}=816\). 若三个编号组成公差为 \(3\) 的等差数列，则它们必为 \(a\)，\(a+3\)，\(a+6\)，其中 \(a=1,2,\ldots,12\)，所以有利事件数为 \(12\). 因此所求概率为 \(\frac{12}{816}=\frac1{68}\)，故选 B.
\end{solution}

\begin{problem}
如图是根据《山东统计年鉴 2007》中的资料作成的 1997 年至 2006 年我省城镇居民百户家庭人口数的茎叶图. 图中左边的数字从左到右分别表示城镇居民百户家庭人口数的百位数字和十位数字，右边的数字表示城镇居民百户家庭人口数的个位数字. 从图中可以得到 1997 年至 2006 年我省城镇居民百户家庭人口数的平均数为（\quad）

\begin{center}
\begin{tabular}{r|l}
\(29\) & \(1\quad1\quad5\quad8\) \\
\(30\) & \(2\quad6\) \\
\(31\) & \(0\quad2\quad4\quad7\)
\end{tabular}
\end{center}

\choices
    {\(304.6\)}
    {\(303.6\)}
    {\(302.6\)}
    {\(301.6\)}
\end{problem}

\begin{answer}
B
\end{answer}

\begin{solution}
由茎叶图可得这十个数据为 \(291\)，\(291\)，\(295\)，\(298\)，\(302\)，\(306\)，\(310\)，\(312\)，\(314\)，\(317\). 它们的和为 \(3036\)，所以平均数为 \(\frac{3036}{10}=303.6\)，故选 B.
\end{solution}

\begin{problem}
\(\left(x - \frac{1}{\sqrt[3]{x}}\right)^{12}\) 展开式中的常数项为（\quad）

\choices
    {\(-1320\)}
    {\(1320\)}
    {\(-220\)}
    {\(220\)}
\end{problem}

\begin{answer}
C
\end{answer}

\begin{solution}
展开式的通项为
\(T_{k+1}=(-1)^k\mathrm C_{12}^{k}x^{12-k}x^{-\frac{k}{3}}=(-1)^k\mathrm C_{12}^{k}x^{12-\frac{4k}{3}}\)，其中\(0\leqslant k\leqslant12\). 要得到常数项，必须有
\(12-\frac{4k}{3}=0\)，解得\(k=9\). 因而常数项的系数为
\((-1)^9\mathrm C_{12}^{9}=-\mathrm C_{12}^{3}=-220\)，故选 C.
\end{solution}

\begin{problem}
设椭圆 \(C_1\) 的离心率为 \(\frac{5}{13}\)，焦点在 \(x\) 轴上且长轴长为 25. 若曲线 \(C_2\) 上的点到椭圆 \(C_1\) 的两个焦点的距离的差的绝对值等于 8，则曲线 \(C_2\) 的标准方程为（\quad）

\choices
    {\(\frac{x^2}{4^2} - \frac{y^2}{3^2} = 1\)}
    {\(\frac{x^2}{13^2} - \frac{y^2}{5^2} = 1\)}
    {\(\frac{x^2}{3^2} - \frac{y^2}{4^2} = 1\)}
    {\(\frac{x^2}{13^2} - \frac{y^2}{12^2} = 1\)}
\end{problem}

\begin{solution}
按原题给出的长轴长 \(25\)，椭圆 \(C_1\) 的长半轴为 \(a_1=\frac{25}{2}\)，所以其焦距参数为 \(c_1=a_1e=\frac{25}{2}\cdot\frac{5}{13}=\frac{125}{26}\). 曲线 \(C_2\) 上的点到这两个焦点的距离差的绝对值为 \(8\)，因此 \(C_2\) 应为以这两个焦点为焦点的双曲线，且其长半轴满足 \(2a_2=8\)，即 \(a_2=4\). 于是 \(b_2^2=c_1^2-a_2^2=\left(\frac{125}{26}\right)^2-16=\frac{4809}{676}\)，按原题数据得到的双曲线方程应为 \(\frac{x^2}{16}-\frac{y^2}{4809/676}=1\).

但四个选项中只有选项 A 的长半轴为 \(4\)，且该选项的焦距参数为 \(\sqrt{4^2+3^2}=5\)，与 \(c_1=\frac{125}{26}\) 不同；其他选项的长半轴或焦点方向也不符合. 因而原题的“长轴长为 \(25\)”与所给四个选项不相容，按原件无法确定一个正确选项.
\end{solution}

\begin{problem}
已知圆的方程为 \(x^{2} + y^{2} - 6x - 8y = 0\). 设该圆过点 \((3,5)\) 的最长弦和最短弦分别为 \(AC\) 和 \(BD\)，则四边形 \(ABCD\) 的面积为（\quad）

\choices
    {\(10\sqrt{6}\)}
    {\(20\sqrt{6}\)}
    {\(30\sqrt{6}\)}
    {\(40\sqrt{6}\)}
\end{problem}

\begin{answer}
B
\end{answer}

\begin{solution}
将圆的方程配方得 \((x-3)^2+(y-4)^2=25\)，所以圆心为 \(O=(3,4)\)，半径为 \(5\). 记题中给定的点为 \(P=(3,5)\)，则 \(OP=1\). 过 \(P\) 的最长弦是过 \(P\) 和圆心的直径，因此 \(AC=2\cdot5=10\). 过 \(P\) 的最短弦垂直于 \(OP\)，其半弦长为 \(\sqrt{5^2-1^2}=2\sqrt6\)，所以 \(BD=4\sqrt6\). 四边形 \(ABCD\) 的两条对角线 \(AC\)，\(BD\) 互相垂直，故其面积为 \(\frac12 AC\cdot BD=\frac12\cdot10\cdot4\sqrt6=20\sqrt6\)，故选 B.
\end{solution}

\begin{problem}
设二元一次不等式组 \(\begin{cases}
x + 2y - 19 \geqslant 0,\\
x - y + 8 \geqslant 0,\\
2x + y - 14 \leqslant 0
\end{cases}\) 所表示的平面区域为 \(M\)，使函数 \(y = a^x\)（\(a > 0, a \neq 1\)）的图象过区域 \(M\) 的 \(a\) 的取值范围是（\quad）

\choices
    {\([1,3]\)}
    {\([2,\sqrt{10}]\)}
    {\([2,9]\)}
    {\([\sqrt{10},9]\)}
\end{problem}

\begin{answer}
C
\end{answer}

\begin{solution}
不等式组所表示的区域由三条直线围成. 三条边界直线两两相交所得的三个顶点为
\((1,9)\)，\((2,10)\)，\((3,8)\)，所以区域 \(M\) 是这三个点组成的三角形. 对于 \((x,y)\in M\)，有 \(1\leqslant x\leqslant3\) 且 \(y>0\). 若 \(y=a^x\)，则
\(\ln a=\frac{\ln y}{x}\).

固定 \(x\) 时，\(\frac{\ln y}{x}\) 随 \(y\) 增大而增大. 区域的下边界为 \(y=\frac{19-x}{2}\)，上边界在 \([1,2]\) 上为 \(y=x+8\)，在 \([2,3]\) 上为 \(y=14-2x\). 对下边界，函数
\(\frac{\ln((19-x)/2)}{x}\) 的导数为
\(\frac{-x/(19-x)-\ln((19-x)/2)}{x^2}<0\)，故其最大值和最小值分别在 \(x=1\) 和 \(x=3\) 处取得，分别为 \(\ln9\) 和 \(\ln2\). 对上边界，在 \([1,2]\) 上
\(\frac{\ln(x+8)}{x}\) 的导数为 \(\frac{x/(x+8)-\ln(x+8)}{x^2}<0\)，在 \([2,3]\) 上
\(\frac{\ln(14-2x)}{x}\) 的导数为 \(\frac{-2x/(14-2x)-\ln(14-2x)}{x^2}<0\)，所以整个上边界上的最大值为 \(\ln9\)，最小值为 \(\ln2\). 因而
\(\ln2\leqslant\ln a\leqslant\ln9\)，即 \(2\leqslant a\leqslant9\). 区域 \(M\) 连通，且 \(y^{1/x}\) 连续，所以其中间的每个值都能取得，故选 C.
\end{solution}

\section{填空题}

\begin{problem}
执行下面的程序框图，若 \(p = 0.8\)，则输出的 \(n =\)\fillinblank{}.

\bitmapfigure[max width=6.0cm]{img/2008/shandong_science/q13_fig1.png}
\end{problem}

\begin{answer}
\(4\)
\end{answer}

\begin{solution}
程序开始时 \(n=1\)，\(S=0\). 由于 \(0<0.8\)，第一次循环后 \(S=\frac12\)，\(n=2\)；由于 \(\frac12<0.8\)，第二次循环后 \(S=\frac12+\frac14=\frac34\)，\(n=3\)；由于 \(\frac34<0.8\)，第三次循环后 \(S=\frac34+\frac18=\frac78\)，\(n=4\). 此时 \(\frac78=0.875>0.8\)，循环结束，故输出 \(n=4\).
\end{solution}

\begin{problem}
设函数 \(f(x) = ax^2 + c\)（\(a \neq 0\)）. 若 \(\int_0^1 f(x)\mathrm{d}x = f(x_0)\)，\(0 \leqslant x_0 \leqslant 1\)，则 \(x_0\) 的值为\fillinblank{}.
\end{problem}

\begin{answer}
\(\frac{\sqrt{3}}{3}\)
\end{answer}

\begin{solution}
由题意，
\(\int_0^1f(x)\mathrm{d}x=\int_0^1(ax^2+c)\mathrm{d}x=\frac{a}{3}+c\). 又 \(f(x_0)=ax_0^2+c\)，所以
\(ax_0^2+c=\frac{a}{3}+c\). 由于 \(a\neq0\)，可得 \(x_0^2=\frac13\). 结合 \(0\leqslant x_0\leqslant1\)，得 \(x_0=\frac{\sqrt3}{3}\).
\end{solution}

\begin{problem}
已知 \(a\)，\(b\)，\(c\) 为 \(\triangle ABC\) 的三个内角 \(A\)，\(B\)，\(C\) 的对边，向量 \(\bs{m} = (\sqrt{3}, -1)\)，\(\bs{n} = (\cos A, \sin A)\). 若 \(\bs{m} \perp \bs{n}\)，且 \(a \cos B + b \cos A = c \sin C\)，则角 \(B =\)\fillinblank{}.
\end{problem}

\begin{answer}
\(\frac{\pi}{6}\)
\end{answer}

\begin{solution}
由 \(\bs{m}\perp\bs{n}\)，得
\(\sqrt3\cos A-\sin A=0\)，即 \(\sin A=\sqrt3\cos A\). 因为 \(0<A<\pi\) 且 \(\sin A>0\)，所以 \(\cos A>0\)，从而 \(A=\frac\pi3\).

由余弦定理，
\(\cos B=\frac{a^2+c^2-b^2}{2ac}\)，\(\cos A=\frac{b^2+c^2-a^2}{2bc}\). 因而
\(a\cos B+b\cos A=\frac{a^2+c^2-b^2}{2c}+\frac{b^2+c^2-a^2}{2c}=c\). 题设还给出 \(a\cos B+b\cos A=c\sin C\)，故 \(c=c\sin C\). 由于 \(c>0\)，所以 \(\sin C=1\)，即 \(C=\frac\pi2\). 因此
\(B=\pi-A-C=\pi-\frac\pi3-\frac\pi2=\frac\pi6\).
\end{solution}

\begin{problem}
若不等式 \(|3x - b| < 4\) 的解集中的整数有且仅有 \(1\)，\(2\)，\(3\)，则 \(b\) 的取值范围为\fillinblank{}.
\end{problem}

\begin{answer}
\((5,7)\)
\end{answer}

\begin{solution}
不等式等价于
\(\frac{b-4}{3}<x<\frac{b+4}{3}\). 为使 \(1\)，\(2\)，\(3\) 都在解集中，需要 \(b<7\) 且 \(b>5\). 同时，为使整数 0 不在解集中，需要 \(\frac{b-4}{3}\geqslant0\)，为使整数 4 不在解集中，需要 \(\frac{b+4}{3}\leqslant4\). 后两个条件在 \(5<b<7\) 时自动满足，因此
\(5<b<7\)，即 \(b\in(5,7)\).
\end{solution}

\section{解答题}

\begin{problem}
已知函数 \(f(x) = \sqrt{3}\sin(\omega x + \varphi) - \cos(\omega x + \varphi)\)（\(0 < \varphi < \pi, \omega > 0\)）为偶函数，且函数 \(y = f(x)\) 图象的两相邻对称轴间的距离为 \(\frac{\pi}{2}\).
\begin{enumerate}
\item 求 \(f\left(\frac{\pi}{8}\right)\) 的值；
\item 将函数 \(y = f(x)\) 的图象向右平移 \(\frac{\pi}{6}\) 个单位后，再将得到的图象上各点的横坐标伸长到原来的 4 倍，纵坐标不变，得到函数 \(y = g(x)\) 的图象，求 \(g(x)\) 的单调递减区间.
\end{enumerate}
\end{problem}

\begin{solution}
\begin{enumerate}
\item 将函数化为正弦型，得
\(f(x)=2\sin\left(\omega x+\varphi-\frac\pi6\right)\). 该函数相邻对称轴间的距离为 \(\frac\pi\omega\)，所以 \(\frac\pi\omega=\frac\pi2\)，即 \(\omega=2\). 令 \(\delta=\varphi-\frac\pi6\)，则
\(f(x)=2\sin(2x+\delta)=2\sin2x\cos\delta+2\cos2x\sin\delta\). 因为 \(f(x)\) 是偶函数，\(\sin2x\) 的系数必须为零，故 \(\cos\delta=0\). 由 \(0<\varphi<\pi\)，得 \(-\frac\pi6<\delta<\frac{5\pi}{6}\)，所以 \(\delta=\frac\pi2\)，从而 \(\varphi=\frac{2\pi}{3}\). 因此 \(f(x)=2\cos2x\)，于是
\(f\left(\frac\pi8\right)=2\cos\frac\pi4=\sqrt2\).

\item 先向右平移 \(\frac\pi6\)，得到 \(y=f\left(x-\frac\pi6\right)=2\cos\left(2x-\frac\pi3\right)\). 再将横坐标伸长到原来的 4 倍，得到
\(g(x)=2\cos\left(\frac{x}{2}-\frac\pi3\right)\). 因而
\(g'(x)=-\sin\left(\frac{x}{2}-\frac\pi3\right)\). 当
\(\frac{x}{2}-\frac\pi3\in[2k\pi,(2k+1)\pi]\)，即
\(x\in\left[4k\pi+\frac{2\pi}{3},4k\pi+\frac{8\pi}{3}\right]\)（\(k\in\Z\)）时，\(g'(x)\leqslant0\)，且除端点外严格小于零. 所以 \(g(x)\) 的单调递减区间为
\(\left[4k\pi+\frac{2\pi}{3},4k\pi+\frac{8\pi}{3}\right]\)（\(k\in\Z\)）.
\end{enumerate}
\end{solution}

\begin{problem}
甲，乙两队参加奥运知识竞赛，每队 \(3\) 人，每人回答一个问题. 答对者为本队赢得一分，答错得零分. 假设甲队中每人答对的概率均为 \(\frac{2}{3}\)，乙队中 \(3\) 人答对的概率分别为 \(\frac{2}{3}\)，\(\frac{2}{3}\)，\(\frac{1}{2}\) 且各人正确与否相互之间没有影响. 用 \(\xi\) 表示甲队的总得分.
\begin{enumerate}
\item 求随机变量 \(\xi\) 分布列和数学期望；
\item 用 \(A\) 表示“甲，乙两个队总得分之和等于 \(3\)”这一事件，用 \(B\) 表示“甲队总得分大于乙队总得分”这一事件，求 \(P(AB)\).
\end{enumerate}
\end{problem}

\begin{solution}
\begin{enumerate}
\item 甲队三人的答题结果相互独立且答对概率均为 \(\frac23\)，所以 \(\xi\) 服从二项分布 \(B\left(3,\frac23\right)\). 因而
\(P(\xi=0)=\left(\frac13\right)^3=\frac1{27}\)，
\(P(\xi=1)=\mathrm C_3^1\frac23\left(\frac13\right)^2=\frac2{9}\)，
\(P(\xi=2)=\mathrm C_3^2\left(\frac23\right)^2\frac13=\frac4{9}\)，
\(P(\xi=3)=\left(\frac23\right)^3=\frac8{27}\). 这些概率之和为 1，故这就是 \(\xi\) 的分布列. 数学期望为
\(E(\xi)=0\cdot\frac1{27}+1\cdot\frac2{9}+2\cdot\frac4{9}+3\cdot\frac8{27}=2\).

\item 记乙队总得分为 \(\eta\). 乙队三人的答题结果相互独立，所以
\(P(\eta=0)=\frac13\cdot\frac13\cdot\frac12=\frac1{18}\)，
\(P(\eta=1)=\frac23\cdot\frac13\cdot\frac12+\frac13\cdot\frac23\cdot\frac12+\frac13\cdot\frac13\cdot\frac12=\frac5{18}\)，
\(P(\eta=2)=\frac23\cdot\frac23\cdot\frac12+\frac23\cdot\frac13\cdot\frac12+\frac13\cdot\frac23\cdot\frac12=\frac4{9}\)，
\(P(\eta=3)=\frac23\cdot\frac23\cdot\frac12=\frac2{9}\).
\(A\cap B\) 只能发生在 \((\xi,\eta)=(2,1)\) 或 \((3,0)\). 由于两队答题结果相互独立，
\(P(AB)=P(\xi=2)P(\eta=1)+P(\xi=3)P(\eta=0)=\frac4{9}\cdot\frac5{18}+\frac8{27}\cdot\frac1{18}=\frac{34}{243}\).
\end{enumerate}
\end{solution}

\begin{problem}
将数列 \(\{a_n\}\) 中的所有项按每一行比上一行多一项的规则排成如下数表：

\begin{center}
\begin{tabular}{cccc}
 &  & \(a_1\) &  \\
 & \(a_2\) & \(a_3\) &  \\
\(a_4\) & \(a_5\) & \(a_6\) &  \\
\(a_7\) & \(a_8\) & \(a_9\) & \(a_{10} \cdots\)
\end{tabular}
\end{center}

记表中的第一列数 \(a_1\)，\(a_2\)，\(a_4\)，\(a_7\)，\(\cdots\) 构成的数列为 \(\{b_n\}\)，\(b_1 = a_1 = 1\). \(S_n\) 为数列 \(\{b_n\}\) 的前 \(n\) 项和，且满足 \(\frac{2b_n}{b_nS_n - S_n^2} = 1\)（\(n \geqslant 2\)）.
\begin{enumerate}
\item 证明数列 \(\left\{\frac{1}{S_n}\right\}\) 成等差数列，并求数列 \(\{b_n\}\) 的通项公式；
\item 上表中，若从第三行起，每一行中的数按从左到右的顺序均构成等比数列，且公比为同一个正数. 当 \(a_{81} = -\frac{4}{91}\) 时，求上表中第 \(k\) 行（\(k \geqslant 3\)）所有项和的和.
\end{enumerate}
\end{problem}

\begin{solution}
\begin{enumerate}
\item 当 \(n\geqslant2\) 时，\(b_n=S_n-S_{n-1}\)，所以
\(b_nS_n-S_n^2=S_n(b_n-S_n)=-S_nS_{n-1}\). 由题设条件，
\(\frac{2b_n}{-S_nS_{n-1}}=1\)，即 \(2b_n=-S_nS_{n-1}\). 代入 \(b_n=S_n-S_{n-1}\)，得
\(2(S_n-S_{n-1})=-S_nS_{n-1}\)，从而 \(S_n(S_{n-1}+2)=2S_{n-1}\). 因为 \(S_1=b_1=1\)，所以
\(\frac1{S_n}=\frac{S_{n-1}+2}{2S_{n-1}}=\frac1{S_{n-1}}+\frac12\). 因此 \(\left\{\frac1{S_n}\right\}\) 是首项为 1、公差为 \(\frac12\) 的等差数列，故
\(\frac1{S_n}=1+\frac{n-1}{2}=\frac{n+1}{2}\)，即 \(S_n=\frac2{n+1}\). 于是
\(b_1=1\)，且当 \(n\geqslant2\) 时
\(b_n=S_n-S_{n-1}=\frac2{n+1}-\frac2n=-\frac2{n(n+1)}\).

\item 第 \(k\) 行的第一项是 \(b_k\)，且第 \(k\) 行有 \(k\) 项. 由于 \(a_{81}\) 位于第 13 行的第 3 列，且 \(b_{13}=a_{79}=-\frac1{91}\)，设各行公比为 \(q>0\)，则
\(a_{81}=b_{13}q^2\)，从而 \(-\frac1{91}q^2=-\frac4{91}\)，得 \(q=2\). 因此第 \(k\) 行所有项的和为
\(b_k(1+2+\cdots+2^{k-1})=-\frac2{k(k+1)}(2^k-1)=\frac{2(1-2^k)}{k(k+1)}\)，其中 \(k\geqslant3\).
\end{enumerate}
\end{solution}

\begin{problem}
如图，已知四棱锥 \(P - ABCD\)，底面 \(ABCD\) 为菱形，\(PA \perp\) 平面 \(ABCD\)，\(\angle ABC = 60^{\circ}\)，\(E\)，\(F\) 分别是 \(BC\)，\(PC\) 的中点.
\begin{enumerate}
\item 证明：\(AE \perp PD\)；
\item 若 \(H\) 为 \(PD\) 上的动点，\(EH\) 与平面 \(PAD\) 所成最大角的正切值为 \(\frac{\sqrt{6}}{2}\)，求二面角 \(E - AF - C\) 的余弦值.
\end{enumerate}

\bitmapfigure[max width=0.32\linewidth]{img/2008/shandong_science/q20_fig1.png}
\end{problem}

\begin{solution}
\begin{enumerate}
\item 设菱形边长为 \(s>0\)，设 \(PA=h>0\). 建立直角坐标系，使
\(A=(0,0,0)\)，\(B=(s,0,0)\)，\(C=\left(\frac{s}{2},\frac{\sqrt3s}{2},0\right)\)，\(D=\left(-\frac{s}{2},\frac{\sqrt3s}{2},0\right)\)，\(P=(0,0,h)\). 则
\(E=\left(\frac{3s}{4},\frac{\sqrt3s}{4},0\right)\)，从而
\(\overrightarrow{AE}=\left(\frac{3s}{4},\frac{\sqrt3s}{4},0\right)\)，
\(\overrightarrow{PD}=\left(-\frac{s}{2},\frac{\sqrt3s}{2},-h\right)\). 因此
\(\overrightarrow{AE}\cdot\overrightarrow{PD}=-\frac{3s^2}{8}+\frac{3s^2}{8}=0\)，故 \(AE\perp PD\).

\item 取 \(H=P+t(D-P)\)，其中 \(0\leqslant t\leqslant1\). 平面 \(PAD\) 的一个单位法向量为 \(\bs{n}=\left(\frac{\sqrt3}{2},\frac12,0\right)\). 令 \(\bs{v}=\overrightarrow{EH}\)，直接计算得
\(\bs{n}\cdot\bs{v}=-\frac{\sqrt3s}{2}\)，且
\(\lvert\bs{v}\rvert^2-\left(\bs{n}\cdot\bs{v}\right)^2=s^2t^2+h^2(1-t)^2\). 设 \(EH\) 与平面 \(PAD\) 所成的角为 \(\theta\)，则
\(\tan\theta=\frac{\sqrt3s/2}{\sqrt{s^2t^2+h^2(1-t)^2}}\). 又
\(s^2t^2+h^2(1-t)^2=(s^2+h^2)\left(t-\frac{h^2}{s^2+h^2}\right)^2+\frac{s^2h^2}{s^2+h^2}\)，所以当 \(t=\frac{h^2}{s^2+h^2}\) 时，\(\tan\theta\) 取得最大值
\(\frac{\sqrt3\sqrt{s^2+h^2}}{2h}\). 由题设最大值为 \(\frac{\sqrt6}{2}\)，得
\(\frac{3(s^2+h^2)}{4h^2}=\frac64\)，即 \(s^2=h^2\). 因此 \(s=h\).

\(s=h\) 时可令 \(s=h=1\) 计算二面角，因为整体缩放不改变角度. 此时
\(A=(0,0,0)\)，\(C=\left(\frac12,\frac{\sqrt3}{2},0\right)\)，
\(E=\left(\frac34,\frac{\sqrt3}{4},0\right)\)，\(F=\left(\frac14,\frac{\sqrt3}{4},\frac12\right)\). 取平面 \(EAF\) 和 \(AFC\) 的法向量
\(\bs{n}_1=\overrightarrow{AE}\times\overrightarrow{AF}\) 与 \(\bs{n}_2=\overrightarrow{AF}\times\overrightarrow{AC}\)，去掉不影响夹角的正比例因子后可取
\(\bs{n}_1=(\sqrt3,-3,\sqrt3)\)，\(\bs{n}_2=(-\sqrt3,1,0)\). 所以二面角 \(E-AF-C\) 的余弦值为
\(\frac{|\bs{n}_1\cdot\bs{n}_2|}{|\bs{n}_1||\bs{n}_2|}=\frac6{2\sqrt{15}}=\frac{\sqrt{15}}5\).
\end{enumerate}
\end{solution}

\begin{problem}
已知函数 \(f(x) = \frac{1}{(1 - x)^n} + a \ln(x - 1)\)，其中 \(n \in \mathbf{N}^*\)，\(a\) 为常数.
\begin{enumerate}
\item 当 \(n = 2\) 时，求函数 \(f(x)\) 的极值.
\item 当 \(a = 1\) 时，证明：对任意的正整数 \(n\)，当 \(x \geqslant 2\) 时，有 \(f(x) \leqslant x - 1\).
\end{enumerate}
\end{problem}

\begin{solution}
\begin{enumerate}
\item 当 \(n=2\) 时，定义域为 \(x>1\)，且令 \(t=x-1>0\)，则
\(f(x)=\frac1{t^2}+a\ln t\)，
\(f'(x)=\frac{at^2-2}{t^3}\). 当 \(a\leqslant0\) 时，\(f'(x)<0\)，所以函数在 \((1,+\infty)\) 上单调递减，没有极值. 当 \(a>0\) 时，令 \(f'(x)=0\)，得 \(t=\sqrt{\frac2a}\)，即 \(x=1+\sqrt{\frac2a}\). 在该点左侧 \(f'(x)<0\)，右侧 \(f'(x)>0\)，所以该点取得极小值，极小值为
\(\frac{a}{2}+a\ln\sqrt{\frac2a}=\frac a2\left(1+\ln\frac2a\right)\). 因而当 \(a>0\) 时，函数的极小值为 \(\frac a2\left(1+\ln\frac2a\right)\)，无极大值.

\item 当 \(a=1\) 且 \(x\geqslant2\) 时，令 \(t=x-1\geqslant1\). 注意到 \(\frac{1}{(1-x)^n}=(-1)^nt^{-n}\). 当 \(n\) 为偶数时，\((-1)^nt^{-n}=t^{-n}\leqslant1\)；当 \(n\) 为奇数时，\((-1)^nt^{-n}=-t^{-n}\leqslant1\)，所以对任意正整数 \(n\)，均有 \(\frac{1}{(1-x)^n}\leqslant1\). 另一方面，函数 \(t-1-\ln t\) 在 \([1,+\infty)\) 上的导数为 \(1-\frac1t\geqslant0\)，且在 \(t=1\) 时为 0，故 \(\ln t\leqslant t-1\). 从而 \(f(x)=(-1)^nt^{-n}+\ln t\leqslant1+(t-1)=t=x-1\).
\end{enumerate}
\end{solution}

\begin{problem}
如图，设抛物线方程为 \(x^{2} = 2py\)（\(p > 0\)），\(M\) 为直线 \(y = -2p\) 上任意一点，过 \(M\) 引抛物线的切线，切点分别为 \(A\)，\(B\).
\begin{enumerate}
\item 求证：\(A\)，\(M\)，\(B\) 三点的横坐标成等差数列；
\item 已知当 \(M\) 点的坐标为 \((2, -2p)\) 时，\(|AB| = 4\sqrt{10}\)，求此时抛物线的方程；
\item 是否存在点 \(M\)，使得点 \(C\) 关于直线 \(AB\) 的对称点 \(D\) 在抛物线 \(x^2 = 2py\)（\(p > 0\)）上，其中，点 \(C\) 满足 \(\overrightarrow{OC} = \overrightarrow{OA} + \overrightarrow{OB}\)（\(O\) 为坐标原点）. 若存在，求出所有适合题意的点 \(M\) 的坐标；若不存在，请说明理由.
\end{enumerate}

\bitmapfigure[max width=0.32\linewidth]{img/2008/shandong_science/q22_fig1.png}
\end{problem}

\begin{solution}
\begin{enumerate}
\item 设切点 \(A\)，\(B\) 的横坐标分别为 \(u\)，\(v\)，则抛物线可写为 \(y=\frac{x^2}{2p}\). 在横坐标为 \(t\) 的点处，切线方程为
\(y=\frac{t}{p}x-\frac{t^2}{2p}\). 设 \(M=(m,-2p)\)，把 \(M\) 代入两条切线，得到 \(u\)，\(v\) 是方程
\(t^2-2mt-4p^2=0\) 的两个根. 由韦达定理，\(u+v=2m\)，所以 \(m=\frac{u+v}{2}\)，即 \(A\)，\(M\)，\(B\) 三点的横坐标成等差数列.

\item 当 \(M=(2,-2p)\) 时，\(u+v=4\)，且 \(uv=-4p^2\). 因此
\(u-v=\sqrt{(u+v)^2-4uv}=4\sqrt{1+p^2}\). 又
\(y_A-y_B=\frac{u^2-v^2}{2p}=\frac{(u-v)(u+v)}{2p}=\frac{2(u-v)}p\). 从而
\(\lvert AB\rvert=\sqrt{(u-v)^2+(y_A-y_B)^2}=\frac4p\sqrt{(1+p^2)(p^2+4)}\). 由 \(\lvert AB\rvert=4\sqrt{10}\)，得
\(\frac{(1+p^2)(p^2+4)}{p^2}=10\). 令 \(q=p^2>0\)，则 \((q+1)(q+4)=10q\)，即 \(q^2-5q+4=0\)，所以 \(q=1\) 或 \(q=4\). 因而 \(p=1\) 或 \(p=2\)，抛物线方程为 \(x^2=2y\) 或 \(x^2=4y\).

\item 设 \(M=(m,-2p)\). 由上面的根关系，\(u+v=2m\)，\(uv=-4p^2\). 因而
\(C=A+B=\left(u+v,\frac{u^2+v^2}{2p}\right)=\left(2m,\frac{2m^2}{p}+4p\right)\).

两切点所在的弦 \(AB\) 的斜率为 \(\frac{u+v}{2p}=\frac{m}{p}\)，其方程为 \(y=\frac{m}{p}x+2p\)，即 \(-mx+py-2p^2=0\). 由点关于直线的对称公式，点 \(C\) 关于此直线的对称点 \(D=(x_D,y_D)\) 为
\(D=C-\frac{4p^2}{m^2+p^2}(-m,p)=\left(\frac{2m(m^2+3p^2)}{m^2+p^2},\frac{2m^2(m^2+3p^2)}{p(m^2+p^2)}\right)\).

由此 \(y_D=\frac{m}{p}x_D\). 若 \(D\) 在抛物线 \(x^2=2py\) 上，则
\(x_D^2=2py_D=2mx_D\)，即 \(x_D(x_D-2m)=0\). 代入 \(x_D=\frac{2m(m^2+3p^2)}{m^2+p^2}\)，若 \(m\neq0\)，则 \(x_D=0\) 不可能，且 \(x_D=2m\) 将推出 \(m^2+3p^2=m^2+p^2\)，与 \(p>0\) 矛盾. 所以只能有 \(m=0\). 当 \(m=0\) 时，\(D=(0,0)\) 正好是抛物线的顶点，故存在且所有符合条件的点只有 \(M=(0,-2p)\).
\end{enumerate}
\end{solution}
